% © 2025--2026 Ing. David Jaroš — CC BY-NC-ND 4.0
\documentclass[11pt,a4paper]{article}
\usepackage{amsmath,amssymb,amsthm,mathrsfs}
\usepackage[most]{tcolorbox}
\usepackage{hyperref}
\usepackage[margin=2.5cm]{geometry}

\newtheorem{definition}{Definition}
\newtheorem{proposition}{Proposition}
\newtheorem{theorem}{Theorem}
\newtheorem{remark}{Remark}

\title{Canonical Metric Definition\\
\large Complex-time-fiber completion of the UBT metric channel}
\author{Ing. David Jaroš}
\date{14 July 2026}

\begin{document}
\maketitle

\begin{abstract}
The UBT metric is derived from the existing biquaternion-valued field
$\Theta(q,\tau)$ and is not independently varied.  This canonical revision
separates the full biquaternionic quadratic product from its scalar metric
projection, removes the local normalization that freezes the lapse, and retains
the complete compact $\psi=\operatorname{Im}\tau$ profile before taking the
classical metric projection.  This fiber-completed definition is the one used by
the rigorous pure-$\Theta$ rank theorem.
\end{abstract}

\section{Fundamental field and raw quadratic tensor}

The core field is
\begin{equation}
 \Theta(q,\tau)\in\mathbb B=\mathbb C\otimes\mathbb H,
 \qquad \dim_{\mathbb R}\mathbb B=8,
 \qquad \tau=x^0+i\psi.
\end{equation}
At fixed $(x,\psi)$ define
\begin{equation}
 \mathcal Q_{\mu\nu}
 :=(\partial_\mu\Theta)(\partial_\nu\Theta)^\ddagger\in\mathbb B.
\end{equation}
This is the full raw biquaternionic quadratic tensor.  Its scalar part
$\operatorname{Sc}(\mathcal Q_{\mu\nu})$ and its ordinary matrix trace are
complex scalars, not full biquaternions.

\section{Fiber pairing}

Let
\begin{equation}
 H(A,B)=\operatorname{Re}\operatorname{Tr}
 \left(A^\ddagger\eta_{\mathbb B}B\right)
\end{equation}
be a fixed real symmetric non-degenerate bilinear on the selected
metric-generating representation sector.  Define
\begin{equation}
 \langle U,V\rangle_\psi
 :=\frac1{2\pi R_\psi}\int_0^{2\pi R_\psi}
 H(U(\psi),V(\psi))\,d\psi.
\end{equation}

\section{Canonical physical metric}

\begin{definition}[Fiber-completed metric]
\begin{equation}
\boxed{
 g_{\mu\nu}[\Theta](x)
 =\frac1{\mathcal N_0}
 \langle\partial_\mu\Theta,\partial_\nu\Theta\rangle_\psi,}
 \qquad \mathcal N_0>0\text{ constant}.
\end{equation}
The constant $\mathcal N_0$ is fixed by asymptotic Minkowski normalization.
\end{definition}

\begin{proposition}[Symmetry and reality]
$g_{\mu\nu}$ is real and symmetric.
\end{proposition}
\begin{proof}
The bilinear $H$ and the fiber integral are real and symmetric.
\end{proof}

\begin{proposition}[Local normalization no-go]
The alternative local normalization
\begin{equation}
 \mathcal N(x)=
 \left|H(\partial_0\Theta,\partial_0\Theta)\right|
\end{equation}
forces $g_{00}=-1$ on a timelike branch and cannot represent a variable lapse.
\end{proposition}
\begin{proof}
It gives $g_{00}=H_{00}/|H_{00}|=-1$ whenever $H_{00}<0$.
\end{proof}

\section{Admissibility}

A field is first-jet admissible when the fiber-averaged Gram matrix is
non-degenerate with Lorentzian inertia:
\begin{equation}
 \det(g_{\mu\nu})\ne0,
 \qquad \operatorname{Inertia}(g)=(1,3)
\end{equation}
for the $(-,+,+,+)$ convention.  Linear independence in a positive-definite
auxiliary norm is not, by itself, a proof of Lorentzian signature.

\section{Restricted section metric}

For diagnostics one may use
\begin{equation}
 g^{(\psi_0)}_{\mu\nu}
 =\mathcal N_0^{-1}
 H(\partial_\mu\Theta(x,\psi_0),\partial_\nu\Theta(x,\psi_0)).
\end{equation}
This section metric is not sufficient for generic pure-$\Theta$ off-shell GR
closure because its value space is only eight-real-dimensional.

\section{Closure status}

The exact first variation and rank theorem are proved in
\begin{center}
\texttt{canonical/gr\_closure/pure\_ubt\_fiber\_closure.tex}.
\end{center}
For the fiber-completed projection, the local kinematic part of GAP-10 is closed
on the explicitly defined fiber-free sector and that condition is open and dense
in local analytic two-jets.  Remaining open items are the one-action separation
of internal and metric equations, the selected Jacobi-sector theorem, and global
closure.

\end{document}
